% Book-specific Chinese adaptation of the scanned AMS FGA Explained interior.
% Original scan: 173.3 x 258.6 mm; centered chapter/section headings,
% unboxed statements, running heads, no decorative common project template.
\documentclass[10pt,twoside,openright,fontset=none,scheme=plain]{ctexbook}
\usepackage[paperwidth=173.3mm,paperheight=258.6mm,inner=17mm,outer=15mm,
  top=20mm,bottom=20mm,headheight=13pt,headsep=7mm,footskip=10mm]{geometry}
\setmainfont{Latin Modern Roman}
\setsansfont{Latin Modern Sans}
\setmonofont{Latin Modern Mono}
\setCJKmainfont{Noto Serif CJK SC}[BoldFont={Noto Serif CJK SC Bold},
 ItalicFont={Noto Serif CJK SC},BoldItalicFont={Noto Serif CJK SC Bold}]
\setCJKsansfont{Noto Sans CJK SC}
\setCJKmonofont{Noto Sans Mono CJK SC}
\usepackage{amsmath,amssymb,amsthm,mathrsfs,bm}
\usepackage{tikz,tikz-cd}
\usetikzlibrary{arrows.meta,calc,patterns,decorations.pathmorphing}
\usepackage{enumitem,longtable,booktabs,array}
\usepackage{fancyhdr}
\usepackage[unicode,hidelinks,bookmarksnumbered]{hyperref}
\hypersetup{pdftitle={基础代数几何：Grothendieck 的 FGA 详解（中文重建工作稿）},
 pdfauthor={Barbara Fantechi; Lothar Göttsche; Luc Illusie; Steven L. Kleiman; Nitin Nitsure; Angelo Vistoli}}
\ctexset{
 part={name={第,部分},number=\arabic{part},format=\centering\Large\bfseries},
 chapter={name={第,章},number=\arabic{chapter},format=\centering\Large\bfseries,
   nameformat=\normalsize\normalfont,titleformat=\Large\bfseries,
   beforeskip=25pt,afterskip=23pt},
 section={format=\centering\normalsize\bfseries,beforeskip=16pt,afterskip=8pt},
 subsection={format=\normalsize\bfseries,beforeskip=11pt,afterskip=5pt}}
\setlength{\parindent}{1.5em}
\setlength{\parskip}{0pt}
\setlength{\emergencystretch}{2em}
\linespread{1.12}
\setlist[enumerate,1]{label=(\roman*),leftmargin=2.5em,itemsep=1pt,topsep=3pt}
\setlist[itemize]{leftmargin=2em,itemsep=1pt,topsep=3pt}
\pagestyle{fancy}
\fancyhf{}
\fancyhead[LE,RO]{\small\thepage}
\fancyhead[CE]{\small\nouppercase{\leftmark}}
\fancyhead[CO]{\small\nouppercase{\rightmark}}
\renewcommand{\headrulewidth}{0pt}
\renewcommand{\chaptermark}[1]{\markboth{\thechapter.\ #1}{}}
\renewcommand{\sectionmark}[1]{\markright{\thesection.\ #1}}
\fancypagestyle{plain}{\fancyhf{}\fancyfoot[C]{\small\thepage}}
\newtheoremstyle{fga}{6pt}{4pt}{\normalfont}{1.5em}{\bfseries}{.}{.5em}{}
\theoremstyle{fga}
\newtheorem{theorem}{定理}[chapter]
\newtheorem{proposition}[theorem]{命题}
\newtheorem{lemma}[theorem]{引理}
\newtheorem{corollary}[theorem]{推论}
\newtheorem{definition}[theorem]{定义}
\newtheorem{remark}[theorem]{注}
\newtheorem{example}[theorem]{例}
\newtheorem{exercise}[theorem]{习题}
\newtheorem{notation}[theorem]{记号}
\newtheorem{convention}[theorem]{约定}
\newenvironment{origstatement}[2]{\par\addvspace{6pt}\noindent\hspace*{1.5em}\textbf{#1 #2.}\ \ignorespaces}{\par\addvspace{4pt}}
\renewcommand{\proofname}{证明}
\numberwithin{equation}{chapter}
\DeclareMathOperator{\Spec}{Spec}
\DeclareMathOperator{\Proj}{Proj}
\DeclareMathOperator{\Hom}{Hom}
\DeclareMathOperator{\End}{End}
\DeclareMathOperator{\Aut}{Aut}
\DeclareMathOperator{\Ob}{Ob}
\DeclareMathOperator{\Mor}{Mor}
\DeclareMathOperator{\Ker}{Ker}
\DeclareMathOperator{\Coker}{Coker}
\DeclareMathOperator{\im}{im}
\DeclareMathOperator{\Pic}{Pic}
\DeclareMathOperator{\Div}{Div}
\DeclareMathOperator{\Hilb}{Hilb}
\DeclareMathOperator{\Quot}{Quot}
\DeclareMathOperator{\Supp}{Supp}
\DeclareMathOperator{\Ext}{Ext}
\DeclareMathOperator{\Tor}{Tor}
\DeclareMathOperator{\Sym}{Sym}
\DeclareMathOperator{\GL}{GL}
\DeclareMathOperator{\Gr}{Gr}
\DeclareMathOperator{\rk}{rk}
\DeclareMathOperator{\length}{length}
\DeclareMathOperator{\trdeg}{trdeg}
\newcommand{\id}{\mathrm{id}}
\newcommand{\op}{\mathrm{op}}
\newcommand{\Sch}{\mathrm{Sch}}
\newcommand{\Set}{\mathrm{Set}}
\newcommand{\cO}{\mathcal O}
\newcommand{\cF}{\mathcal F}
\newcommand{\cG}{\mathcal G}
\newcommand{\cH}{\mathcal H}
\newcommand{\cI}{\mathcal I}
\newcommand{\cL}{\mathcal L}
\newcommand{\cC}{\mathcal C}
\newcommand{\cD}{\mathcal D}
\newcommand{\cE}{\mathcal E}
\newcommand{\defeq}{\mathrel{\mathop:}=}
\newcommand{\pr}{\operatorname{pr}}
\newcommand{\partauthor}[1]{\begin{center}\large #1\end{center}}
